% Gesamttest «Flächen, Pythagoras und Ähnlichkeit» — Quelle des PDFs (python3 scripts/build-lp-pdf.py pythagoras).
% Musterlösung und Raster: bewertungspaket.tex. Alle Werte mit python3 nachgerechnet (scripts/lp/pythagoras/zahlen.py).
% Vorgabe ~/TODO-BM-Vorkurs-LP03-15.md §6: Flächen mit Grundseite und Höhe, zusammengesetzt (G1, G2), Pythagoras (G3, G4),
% Hilfsdreieck und Wasserspiegel im Rohr (G5, G6), Regal mit ähnlichen Dreiecken, k und k², Strahlensatz am Schatten (G7, G8).
% 25 Punkte, rund 30 Minuten. Alle Wurzeln gehen auf: ohne Taschenrechner.
\documentclass[11pt]{article}
\usepackage{../lp-druck}
\renewcommand{\lpname}{Flächen, Pythagoras und Ähnlichkeit}
\renewcommand{\lpteilgebiet}{}
\renewcommand{\lpfuss}{Gesamttest · 25 Punkte}
\begin{document}
\lpkopf{Gesamttest}{%
  \vspace{4pt}\noindent\begin{tabularx}{\linewidth}{@{}X X@{}}
  Code (kein Name): \hrulefill & Punkte: \hrulefill\ / 25
  \end{tabularx}}

\begin{lpinfo}{Anleitung}
\textbf{Zeit:} rund 30 Minuten. \textbf{Hilfsmittel:} Lineal, kein Taschenrechner. Alle Wurzeln gehen auf.
Schreib zu jeder Aufgabe den \textbf{Rechenweg} auf und zeichne, wo verlangt, das \textbf{Hilfsdreieck} ein; bewertet wird der Weg.
Danach: selbst mit dem \emph{Bewertungspaket} bewerten oder die Lösung samt diesem Blatt scannen bzw.\ fotografieren
(ohne Namen und Standort) und mit dem Bewertungspaket einer KI geben.
\end{lpinfo}

\lpteil{Teil A · Grundseite und Höhe}{Kapitel 1 · 6 Punkte}

\begin{lpaufgabe}{G1}{Grundstück}{3 P}
Ein Grundstück hat die Form eines Trapezes (Masse in m, die schrägen Seiten gerundet). Berechne seine Fläche.
\begin{center}
\begin{tikzpicture}[x=4.5mm,y=4.5mm,font=\small]
  \draw[lpblau,line width=1pt,fill=lpblau!8] (0,0) -- (15,0) -- (11,8) -- (2,8) -- cycle;
  \draw[lporange,line width=0.8pt,dashed] (2,0) -- (2,8);
  \draw[line width=0.5pt] (2,0) ++(0.5,0) |- ++(-0.5,0.5);
  \node[below] at (7.5,0) {15}; \node[above] at (6.5,8) {9}; \node[right] at (2,4) {8};
  \node[left] at (1,4) {8.2}; \node[right] at (13,4) {8.9};
\end{tikzpicture}
\end{center}
\lpplatz{4}
\end{lpaufgabe}

\begin{lpaufgabe}{G2}{Dreieck und Parallelogramm}{3 P}
\begin{enumerate}[label=(\alph*),leftmargin=*,itemsep=3pt]
\item Ein Dreieck hat die Fläche $42$~cm$^2$ und die Höhe $7$~cm. Wie lang ist die zugehörige Grundseite?
\item Ein Parallelogramm hat die Grundseite $9$~cm, die Höhe $6$~cm und die schräge Seite $7.5$~cm. Wie gross ist seine Fläche?
\end{enumerate}
\lpplatz{4}
\end{lpaufgabe}

\lpteil{Teil B · Satz des Pythagoras}{Kapitel 2 · 6 Punkte}

\begin{lpaufgabe}{G3}{Pythagoras rechnen und prüfen}{3 P}
\begin{enumerate}[label=(\alph*),leftmargin=*,itemsep=3pt]
\item Ein rechtwinkliges Dreieck hat die Katheten $18$~cm und $24$~cm. Wie lang ist die Hypotenuse?
\item Ist ein Dreieck mit den Seiten $8$~cm, $15$~cm und $18$~cm rechtwinklig? Begründe mit einer Rechnung.
\end{enumerate}
\lpplatz{5}
\end{lpaufgabe}

\begin{lpaufgabe}{G4}{Hypotenuse erkennen}{3 P}
Ein rechtwinkliges Dreieck hat die Seiten $11$~cm, $61$~cm und $x$. Der rechte Winkel liegt zwischen der Seite $11$~cm und der Seite $x$.
Welche Seite ist die Hypotenuse? Berechne $x$.
\lpplatz{5}
\end{lpaufgabe}

\lpteil{Teil C · Hilfsdreiecke und Kreis}{Kapitel 3 · 7 Punkte}

\begin{lpaufgabe}{G5}{Gleichschenkliges Dreieck}{4 P}
Ein gleichschenkliges Dreieck hat die Basis $30$~cm und zwei Schenkel von je $39$~cm.
\begin{enumerate}[label=(\alph*),leftmargin=*,itemsep=3pt]
\item Zeichne eine Skizze und markiere das Hilfsdreieck, mit dem du die Höhe berechnest.
\item Berechne die Höhe.
\item Berechne die Fläche des Dreiecks.
\end{enumerate}
\begin{center}\begin{tikzpicture}[x=5mm,y=5mm]\draw[lplinie,line width=0.3pt,step=1] (0,0) grid (14,6);\end{tikzpicture}\end{center}
\lpplatz{4}
\end{lpaufgabe}

\begin{lpaufgabe}{G6}{Wasser im Rohr}{3 P}
Ein Rohr hat den Innendurchmesser $50$~cm. Das Wasser steht $10$~cm tief (weniger als halb hoch). Die Skizze ist nicht massstäblich.
\begin{center}
\begin{tikzpicture}[x=1.4mm,y=1.4mm,font=\small]
  \fill[lpblau!18] (-16,-12) arc[start angle=216.87,end angle=323.13,radius=20] -- cycle;
  \draw[line width=1pt] (0,0) circle[radius=20];
  \draw[lpblau,line width=1pt] (-16,-12) -- (16,-12);
  \fill (0,0) circle[radius=0.5]; \node[above] at (0,0) {$M$};
  \draw[lporange,<->] (19,-12) -- (19,-20) node[midway,right] {\footnotesize $10$~cm};
\end{tikzpicture}
\end{center}
Wie breit ist der Wasserspiegel? Zeichne das Hilfsdreieck ein.
\lpplatz{5}
\end{lpaufgabe}

\lpteil{Teil D · Ähnliche Figuren}{Kapitel 4 · 6 Punkte}

\begin{lpaufgabe}{G7}{Eckregal}{3 P}
Ein Eckregal hat die Form eines rechtwinkligen Dreiecks mit den Katheten $90$~cm und $60$~cm. Für ein Rohr in der Zimmerecke wird
beim rechten Winkel ein dazu ähnliches Dreieck ausgesägt, dessen längere Kathete $30$~cm lang ist.
\begin{center}
\begin{tikzpicture}[x=0.6mm,y=0.6mm,font=\small]
  \draw[lpblau,line width=1pt,fill=lpblau!8] (0,0) -- (90,0) -- (0,60) -- cycle;
  \draw[lporange,line width=1pt,fill=lporange!20] (0,0) -- (30,0) -- (0,20) -- cycle;
  \draw[|<->|] (0,-17) -- (90,-17) node[midway,below] {$90$~cm};
  \draw[|<->|] (-7,0) -- (-7,60) node[midway,left] {$60$~cm};
  \draw[|<->|,lporange] (0,-5) -- (30,-5) node[midway,below] {\footnotesize $30$~cm};
\end{tikzpicture}
\end{center}
\begin{enumerate}[label=(\alph*),leftmargin=*,itemsep=3pt]
\item Wie gross ist der Streckfaktor $k$ vom grossen zum kleinen Dreieck, und wie lang ist die kürzere Kathete des kleinen Dreiecks?
\item Wie gross ist die Fläche des ausgesägten Dreiecks?
\item Wie gross ist die Restfläche des Regals?
\end{enumerate}
\lpplatz{5}
\end{lpaufgabe}

\begin{lpaufgabe}{G8}{Schatten}{3 P}
Ein $1.2$~m langer Stab wirft einen $0.8$~m langen Schatten. Zur gleichen Zeit wirft ein Turm einen $14$~m langen Schatten.
\begin{enumerate}[label=(\alph*),leftmargin=*,itemsep=3pt]
\item Mach eine Skizze mit den beiden ähnlichen Dreiecken und markiere entsprechende Seiten.
\item Berechne die Höhe des Turms.
\end{enumerate}
\lpplatz{6}
\end{lpaufgabe}
\end{document}
